% VERSION COMPLETE ET VERIFIEE -- 19 juin 2026
% !TEX root = main.tex
% Annexe « Colloques, séminaires et formations »
% À inclure après cartographie.tex : cartographie.tex contient déjà \appendix.

\chapter{Colloques \& formations}
\label{ann:colloques-formations}

\begingroup

% -----------------------------------------------------------------------------
% Charte locale : nuances de gris, dans la continuité de la cartographie.
% -----------------------------------------------------------------------------
\colorlet{ChronoBlack}{black!88}
\colorlet{ChronoDark}{black!70}
\colorlet{ChronoMid}{black!42}
\colorlet{ChronoLight}{black!15}
\colorlet{ChronoPale}{black!5}

% Liens cliquables, mais visuellement fondus dans la charte en niveaux de gris.
\hypersetup{colorlinks=true,allcolors=ChronoBlack,pdfborder={0 0 0}}

\newcommand{\ChronoType}[1]{%
  \begingroup
  \setlength{\fboxsep}{2.5pt}%
  \colorbox{ChronoPale}{%
    \textcolor{ChronoDark}{\scriptsize\bfseries\textsc{#1}}%
  }%
  \endgroup
}

\newcommand{\ChronoAnnee}[1]{%
  \multicolumn{2}{@{}l@{}}{%
    \parbox{\textwidth}{%
      \vspace{.35\baselineskip}%
      {\color{ChronoBlack}\large\bfseries\textsc{#1}}%
      \par\vspace{.22\baselineskip}%
      {\color{ChronoLight}\hrule height .7pt}%
    }%
  }\\*[.8\baselineskip]
}

\newcommand{\ChronoEntree}[3]{%
  {\small\bfseries #1}%
  &
  \ChronoType{#2}\par\vspace{.35\baselineskip}%
  #3%
  \\[1.15\baselineskip]
}

% \begin{center}
%   \setlength{\fboxsep}{9pt}%
%   \colorbox{ChronoPale}{%
%     \parbox{.91\textwidth}{%
%       \small
%       Dans la continuité de la rubrique établie pour le premier CSI, cette
%       chronologie rassemble, par ordre antéchronologique, les principaux
%       colloques, séminaires, écoles thématiques, formations, prix et échanges
%       scientifiques ayant jalonné la thèse. Pour lever toute ambiguïté, la
%       mention \enquote{participation} désigne une présence en tant qu'auditeur
%       ou participant, sans communication scientifique ; les présentations,
%       posters, prix et responsabilités d'organisation sont signalés
%       explicitement. Les intitulés des manifestations, communications et
%       ressources associées sont cliquables.
%     }%
%   }
% \end{center}

\vspace{.7\baselineskip}
\small
\setlength{\LTpre}{0pt}
\setlength{\LTpost}{0pt}
\setlength{\tabcolsep}{0pt}
\renewcommand{\arraystretch}{1.04}

\begin{longtable}{%
  @{}%
  >{\raggedleft\arraybackslash\color{ChronoDark}}p{.15\textwidth}%
  @{\hspace{1.05em}}%
  !{\color{ChronoMid}\vrule width .55pt}%
  @{\hspace{1.20em}}%
  >{\raggedright\arraybackslash\color{black}}p{.76\textwidth}%
  @{}%
}

\endfirsthead
\multicolumn{2}{@{}r@{}}{%
  \scriptsize\color{ChronoMid}\emph{Colloques \& formations --- suite}%
}\\[.65\baselineskip]
\endhead

\ChronoAnnee{2026}

\ChronoEntree{17--20 juin 2026}{Forum technologique}{%
  Participation, sans présentation, avec l'
  \href{https://www.inria.fr/fr/inria-startup-studio}{Inria Startup Studio} à
  \href{https://vivatech.com/}{\emph{VivaTech 2026}}, au parc des exposition de la Porte de Versailles (Paris). Cette présence au plus grand rendez-vous européen des start-ups
  et de la technologie s'inscrit dans la découverte de l'écosystème
  entrepreneurial, des acteurs du transfert et des voies de valorisation de
  résultats issus de la recherche.
}

\ChronoEntree{15 juin 2026}{Séminaire}{%
  Présentation au
  \href{https://team.inria.fr/mathnet/seminar/}{séminaire de l'équipe
  \textsc{MathNet}} du Centre Inria de Paris, intitulée
  \emph{A Bayesian Framework for Community Detection on Graphs}. L'exposé est
  consacré à une formulation bayésienne de la détection de communautés sur les
  graphes, ainsi qu'aux principes d'inférence et d'évaluation qui en découlent.
}

\ChronoEntree{4--7 juin 2026}{Formation entrepreneuriale}{%
  Participation, sans communication scientifique, au programme
  \href{https://newsroom.univ-cotedazur.fr/actualites-evenements/starthese-2026-valorisez-vos-travaux-de-recherche}{%
  \emph{Starthèse 2026}}, organisé à Taglio-Isolaccio (Corse) dans le cadre du
  \href{https://www.pepite-france.fr/starthese-2026-31-dispositifs-pour-accompagner-les-doctorants-vers-lentrepreneuriat/}{%
  challenge national Starthèse}. Ce \emph{bootcamp}
  accompagne les doctorants et jeunes docteurs dans la transformation de
  travaux de recherche en projets entrepreneuriaux et dans le transfert issues de la recherche.
}

\ChronoEntree{3 juin 2026}{Organisation \& entrepreneuriat}{%
  Appui à l'organisation et participation au
  \href{https://www.minesparis.psl.eu/blog/evenements/deeptech-forum-sophia-antipolis-2/}{%
  \emph{DeepTech Forum Sophia Antipolis \#2}}, sur le campus de Mines
  Paris--PSL à Sophia Antipolis. La journée, coorganisée par le mastère
  spécialisé \emph{Deeptech} de Mines Paris--PSL et Dynergie, est consacrée au
  passage du laboratoire au marché, au financement de l'innovation de rupture
  et aux relations entre chercheurs, entrepreneurs, investisseurs et
  industriels. La contribution porte sur l'organisation de l'événement et la
  participation aux échanges. Sans présentation scientifique.
}

\ChronoEntree{7 avril 2026}{Séminaire}{%
  Présentation à l'\href{https://www.ihp.fr/}{Institut Henri \nom{Poincaré}}
  (Paris), dans le cadre du séminaire
  \href{https://imaging-in-paris.github.io/}{\emph{Imaging in Paris}},
  intitulée \emph{From Statistics to Graphs: Hierarchical Clustering for Image
  Analysis}. L'exposé relie la définition statistique des clusters de haute
  densité, leur construction par graphes géométriques et plusieurs applications
  au traitement d'image et de nuages de points LiDAR.
}

\ChronoEntree{17 mars 2026}{Journée scientifique}{%
  Participation en tant qu'auditeur, sans présentation, à la conférence
  \href{https://www.academie-sciences.fr/lintelligence-artificielle-pour-les-sciences}{%
  \emph{L'intelligence artificielle pour les sciences}}, organisée par
  l'Académie des sciences sous la Coupole de l'Institut de France à Paris. La
  journée présente, pour plusieurs disciplines, la manière dont les méthodes
  d'intelligence artificielle complètent les approches mathématiques,
  statistiques et informatiques usuelles.
}

\ChronoEntree{18 janvier 2026}{Publication}{%
  Publication dans la revue
  \href{https://link.springer.com/journal/41109}{\emph{Applied Network Science}}
  de l'article
  \href{https://doi.org/10.1007/s41109-025-00756-1}{\emph{Generalization of
  Single-Linkage with Higher-Order Interactions}} \cite{BibiAppliedNetworkScience}.
  Ce travail développe la généralisation du Single-Linkage fondée sur des
  interactions d'ordre supérieur et en établit les principaux fondements
  théoriques et algorithmiques.
}

\ChronoAnnee{2025}

\ChronoEntree{24 novembre 2025}{Séminaire industriel}{%
  Séminaire au sein de
  \href{https://www.valeo.com/fr/valeo-ai/}{\emph{Valeo AI}} (Paris), intitulé
  \emph{Méthodes de clustering avec graphes. Ou
  \enquote{Comment hacker (H)DBSCAN ?} Théorie, algorithmes \& applications}.
  La présentation propose une lecture unifiée des méthodes de clustering
  hiérarchique, puis montre comment adapter leur critère d'extraction à la
  structure particulière des données étudiées.
}

\ChronoEntree{20 novembre 2025}{Entrepreneuriat \& collaboration}{%
  Participation, sans présentation personnelle, à la
  \href{https://pole-enter.org/evenements/la-fete-des-startups/}{Fête des
  start-ups d'Inria Startup Studio}, organisée au Campus Cyber à Paris-La
  Défense. Cette journée marque le début de la collaboration avec Marie
  \nom{Aspro} autour de la détection non supervisée d'anomalies dans des nuages
  de points LiDAR acquis sur des chantiers et sites industriels.
}

\ChronoEntree{28 octobre 2025}{Conférence}{%
  Présentation en session poster à la
  \href{https://asilomarssc.org/}{\emph{Asilomar Conference on Signals,
  Systems, and Computers 2025}} (Pacific Grove, Californie) de l'article
  \href{https://cmsworkshops.com/Asilomar2025/view_paper.php?PaperNum=1109}{%
  \emph{Higher-Order Monte Carlo Cluster Dynamics for Community Detection in
  Euclidean Graphs}} \cite{BibiMonterey}. Ce travail introduit une dynamique de
  Monte-Carlo fondée sur des interactions triangulaires pour la détection de
  communautés sur des graphes euclidiens frustrés.
}

\ChronoEntree{Octobre 2025}{Séminaire}{%
  Présentation au
  \href{https://3ia.univ-cotedazur.eu/research/doctoral-postdoctoral-seminar-1}{%
  séminaire doctoral et postdoctoral du 3IA Côte d'Azur}
  (\textsc{Université Côte d'Azur}), sous le titre
  \emph{Generalization of Single-Linkage with Higher-Order Interactions}.
}

\ChronoEntree{Juillet puis octobre 2025}{Prix scientifique}{%
  Candidature au
  \href{https://www.cnrs.fr/CNRS-Hebdo/cote-azur/Documents/28/Document.aspx}{%
  prix Pierre \nom{Laffitte}} 2025, sélection parmi les finalistes, puis
  présentation en finale du dossier
  \emph{Classification sur nuages de points euclidiens et graphes aléatoires :
  théorie, algorithmes \& partenariats industriels}. La candidature synthétise
  les résultats théoriques de la thèse ainsi que leurs prolongements avec le
  Cerema et \emph{greehill}.
}

\ChronoEntree{25--29 août 2025}{Conférence}{%
  Présentation au
  \href{https://gretsi.fr/archives/colloque/2025}{\emph{GRETSI 2025}}
  (Strasbourg) de l'article
  \emph{Généralisation du filtre de Frangi pour l'extraction des réseaux de
  fissures au sein d'images d'un glissement de terrain} \cite{BibiStrasbourg}.
  Cette communication est issue de la collaboration engagée avec le Cerema à la
  suite d'EUSIPCO 2024 et compare une approche géométrique non supervisée à une
  méthode d'apprentissage profond.
}

\ChronoEntree{30 juin--12 juillet 2025}{École d'été}{%
  Participation à la 53\ieme{} école d'été de probabilités de Saint-Flour
  (Cantal). Cette école thématique de référence a permis de suivre trois
  cours avancés, dont celui de Justin \nom{Salez} intitulé
  \emph{Modern aspects of Markov chains: entropy, curvature and the cutoff phenomenon}
  \cite{JustinSaintFlour2025}. L'enseignement a notamment porté sur les aspects
  modernes de la théorie des chaînes de Markov, les notions d'entropie et de
  courbure pour l'analyse du phénomène de \emph{cutoff}.
}

\ChronoEntree{11--14 juin 2025}{Formation entrepreneuriale}{%
  Participation, sans présentation, à une \emph{Learning Expedition} organisée
  par \href{https://www.inria.fr/fr/inria-startup-studio}{Inria Startup Studio}
  à l'occasion de
  \href{https://vivatech.com/media/press-releases/viva-tech-explores-the-new-frontiers-of-innovation-11-14-june-2025}{%
  \emph{Viva Technology 2025}} (parc des expositions de la Porte de Versailles, à Paris). Cette
  immersion au grand rendez-vous européen des start-ups et de la technologie
  est consacrée à la découverte de l'écosystème entrepreneurial, à la
  valorisation de résultats de recherche et aux modalités de transfert vers
  l'industrie.
}

\ChronoEntree{10--11 juin 2025}{Workshop scientifique}{%
  Participation en tant qu'auditeur, sans communication, au workshop
  \href{https://indico.math.cnrs.fr/event/14166/}{\emph{Stat-Mech in Créteil
  2025}}, organisé à l'Université Paris-Est Créteil par des membres de l'équipe
  de probabilités du LAMA. Cette édition est consacrée à trois thèmes proches
  des travaux de thèse : la physique statistique, la percolation et les graphes
  aléatoires.
}

\ChronoEntree{24 mars 2025}{Séminaire}{%
  Présentation au
  \href{https://team.inria.fr/mathnet/seminar/}{séminaire de l'équipe
  \textsc{MathNet}/\textsc{Dyogene}} du Centre Inria de Paris, intitulée
  \emph{(Hyper-)graphs, percolation and clustering performances}. L'exposé met
  en regard les performances des méthodes de clustering hiérarchique et la
  vitesse du phénomène de percolation dans les graphes et hypergraphes
  géométriques.
}

\ChronoEntree{13--14 mars 2025}{Formation entrepreneuriale}{%
  Participation, sans présentation, à une \emph{Learning Expedition} de
  \href{https://www.inria.fr/fr/inria-startup-studio}{Inria Startup Studio}
  lors du
  \href{https://hello-tomorrow.org/key-highlights-expert-insights-and-groundbreaking-innovations-in-deep-tech-at-the-10th-edition-of-the-hello-tomorrow-global-summit/}{%
  \emph{Hello Tomorrow Global Summit 2025}} au CENTQUATRE-PARIS, forum
  international consacré à la \emph{deep tech}.
  L'immersion porte sur les modèles de valorisation et sur les rencontres entre
  équipes de recherche, start-ups, industriels et investisseurs.
}

\ChronoEntree{Janvier--mars 2025}{Échanges scientifiques}{%
  Plusieurs discussions techniques avec
  \href{https://www.linkedin.com/in/benjamin-berta/?locale=hu}{Benjamin
  \nom{Berta}}, CTO de
  \href{https://www.greehill.fr/}{\emph{greehill}}, sur la segmentation
  d'instances dans des scènes extérieures acquises par LiDAR, notamment pour
  l'inventaire automatisé d'arbres urbains. Ces échanges prolongent la rencontre
  avec Gyula \nom{Fekete} lors de l'école d'été SSRM 2024 et portent sur
  l'intégration d'une généralisation du Single-Linkage en remplacement de la
  brique DBSCAN du \emph{pipeline} existant. Collaboration à reprendre après la thèse !
}

\ChronoEntree{28--29 janvier 2025}{Colloque}{%
  Participation en tant qu'auditeur, sans communication, au colloque
  \href{https://www.college-de-france.fr/fr/agenda/colloque/geometries-aleatoires-et-applications}{%
  \emph{Géométries aléatoires et applications}}, organisé au Collège de France
  à Paris. Les exposés proposent un panorama des géométries aléatoires, depuis
  les graphes et surfaces aléatoires jusqu'aux applications en imagerie, en
  traitement du signal et dans les télécommunications.
}

\ChronoAnnee{2024}

\ChronoEntree{10--12 décembre 2024}{Conférence}{%
  Présentation à la
  \href{https://2024.complexnetworks.org/}{13\ieme{} conférence internationale
  \emph{Complex Networks and their Applications}} (Istambul) de l'article
  \href{https://doi.org/10.1007/978-3-031-82431-9_2}{%
  \emph{Hypergraphs, Percolation, and Hierarchical Clustering}}
  \cite{BibiIstambul}. Ce travail compare les propriétés de percolation de
  plusieurs généralisations du Single-Linkage et met en évidence l'intérêt des
  interactions d'ordre supérieur.
}

\ChronoEntree{26--30 août 2024}{Conférence}{%
  Présentation en session poster à
  \href{https://eusipcolyon.sciencesconf.org/}{\emph{EUSIPCO 2024}} (Lyon) de
  l'article
  \href{https://doi.org/10.23919/EUSIPCO63174.2024.10715271}{%
  \emph{Benefits of Hypergraphs for Density-Based Clustering}}
  \cite{BibiLyon}; le
  \href{https://overleaf.irisa.fr/read/hnbyhhrcybgb}{poster associé} est
  également disponible. La conférence est aussi à l'origine de la mise en
  relation avec Pierre \nom{Charbonnier} (Cerema), puis de la collaboration sur
  l'extraction automatique de réseaux de fissures.
}

\ChronoEntree{29 juillet--2 août 2024}{École d'été}{%
  Participation à l'école d'été
  \href{https://ssrm.mik.uni-pannon.hu/}{\emph{Remote Sensing and Microscopy
  Image Processing} (SSRM)} à Veszprém (Hongrie), organisée en l'honneur de
  Josiane \nom{Zerubia}. Présentation d'un
  \href{https://overleaf.irisa.fr/read/qtthyrmpjwgs}{poster} et obtention du
  second prix de la
  \href{https://overleaf.irisa.fr/read/wkqjqwhyrtpc}{compétition de recherche
  par équipes}. Une présentation de Gyula \nom{Fekete} sur l'inventaire des
  arbres urbains à partir de LiDAR constitue le point de départ des échanges
  ultérieurs avec \emph{greehill}.
}

\ChronoEntree{24--28 juin 2024}{École d'été}{%
  Participation à l'école d'été
  \href{https://www.normastic.fr/action-transverse-graphes/ecole-dete-graphes/}{%
  \emph{GRAPHADON --- Graphes et leurs applications au traitement de données}}
  à l'INSA Rouen Normandie. La formation couvre notamment la théorie des
  graphes, les graphes sémantiques et dynamiques, le traitement du signal sur
  graphes et l'apprentissage sur données relationnelles.
}

\ChronoEntree{10--14 juin 2024}{Compétition scientifique}{%
  Participation à l'
  \href{https://www.sigmetrics.org/sigmetrics2024/student_activities.html}{%
  ACM Student Research Competition de SIGMETRICS 2024} à Venise avec le travail
  \href{https://doi.org/10.1145/3725536.3725546}{%
  \emph{How Can We Theoretically Measure the Performance of Density-Based
  Clustering Algorithms?}} \cite{BibiVenise}. Après la
  \href{https://overleaf.irisa.fr/read/snxqczzrkqth}{présentation du poster},
  qualification pour la finale orale, puis obtention du deuxième prix dans la
  catégorie des doctorants.
}

\ChronoEntree{3 juin 2024, 14 h}{Séminaire doctoral}{%
  Présentation d'environ trente minutes, suivie d'une discussion, au
  \emph{PhD Seminar} du Centre Inria d'\textsc{Université Côte d'Azur}.
}

\ChronoEntree{27--31 mai 2024}{Colloque}{%
  Participation aux
  \href{https://geosto24.sciencesconf.org/}{Journées de géométrie stochastique}
  à Tours, comprenant des mini-cours et des exposés consacrés aux processus
  ponctuels, à la géométrie aléatoire et à la percolation. Présentation d'un
  \href{https://overleaf.irisa.fr/read/snxqczzrkqth}{poster} sur la mesure
  théorique des performances des algorithmes de clustering fondés sur la
  densité.
}

\ChronoEntree{8 mai 2024}{Workshop}{%
  Exposé à distance au workshop
  \href{https://www.torontomu.ca/graphs-group/hypergraphs2024/}{%
  \emph{Modelling and Mining Complex Networks as Hypergraphs}}, organisé par
  l'Université métropolitaine de Toronto, autour de la généralisation du
  clustering fondé sur des graphes aux interactions d'ordre supérieur.
}

\ChronoEntree{8 avril 2024}{Séminaire invité}{%
  À l'invitation de Gabriele \nom{Moser} et Sebastiano B. \nom{Serpico},
  présentation au département DITEN de l'Université de Gênes du séminaire
  \emph{Density-Based Clustering: (Hyper-)Graphs \& Percolation}.
}

\ChronoEntree{5--6 février 2024}{Colloque}{%
  Participation aux
  \href{https://univ-cotedazur.eu/events/complex-days/past-events/edition-2024}{%
  \emph{Complex Days 2024}} à Nice et présentation des travaux consacrés aux
  graphes géométriques, au clustering hiérarchique et au phénomène de
  percolation.
}

\ChronoAnnee{2023}

\ChronoEntree{28--30 novembre 2023}{Conférence}{%
  Présentation à la
  \href{https://2023.complexnetworks.org/}{12\ieme{} conférence internationale
  \emph{Complex Networks and their Applications}} à Menton du papier
  \href{https://doi.org/10.1007/978-3-031-53472-0_32}{%
  \emph{Graph Based Approach for Galaxy Filament Extraction}}
  \cite{BibiMenton}. Le travail propose un estimateur de densité fondé sur des
  graphes géométriques et l'applique à l'extraction de filaments de galaxies;
  le \href{https://overleaf.irisa.fr/read/qqxqttywnwmv}{poster présenté} est
  accessible en ligne. Le passage aux interactions d'ordre supérieur est évoqué.
}

\ChronoEntree{Octobre 2023}{Formation}{%
  Formation de deux jours aux
  \href{https://www.pssmfrance.fr/}{Premiers secours en santé mentale (PSSM
  France)}, consacrée au repérage des troubles psychiques, à l'écoute initiale
  et à l'orientation vers les dispositifs d'aide adaptés.
}

\ChronoEntree{10--11 octobre 2023}{Mini-cours}{%
  Mini-cours de Dieter \nom{Mitsche} au Centre de Recerca Matemàtica de
  Barcelone,
  \href{https://www.crm.cat/mini-course-on-different-random-geometric-graph-models/}{%
  \emph{Geometric Random Graph Models}}, composé de quatre séances d'une heure
  trente sur les graphes géométriques aléatoires euclidiens et hyperboliques,
  les modèles spatiaux et leurs liens avec la percolation. Une
  \href{https://www.overleaf.com/read/nxvjcnnmqcwv}{retranscription des notes
  de cours} est disponible.
}

\ChronoEntree{7--8 septembre 2023}{Colloque}{%
  Participation aux
  \href{https://geostat23.sciencesconf.org/}{XVI\ieme{} Journées de
  Géostatistique} à Fontainebleau et présentation d'un
  \href{https://overleaf.irisa.fr/read/pnzpyhsxbycf}{poster} consacré à une
  approche par graphes pour l'extraction de filaments de galaxies.
}

\ChronoEntree{28 août--1\ier{} septembre 2023}{École d'été}{%
  Participation à l'école d'été
  \href{https://indico.math.cnrs.fr/event/9939/}{\emph{Geometry and Data}} à
  Strasbourg, consacrée aux outils géométriques pour l'analyse de données, avec
  des cours et travaux pratiques en transport optimal, analyse de formes,
  topologie des données astrophysiques et traitement géométrique des images.
}

\ChronoEntree{Juin 2023}{Formation}{%
  Participation au \emph{Scientific English CLASSROOM Workshop} du Centre Inria
  d'Université Côte d'Azur, formation intensive à la communication scientifique
  en anglais, à l'écrit comme à l'oral.
}

\end{longtable}

\endgroup
